\documentclass[11pt,letterpaper]{article}
\usepackage{technical-user-manual}
\setDocTitle{CI Principles and Integration System Model}
\setDocSubtitle{Engineering technical reference}
\setDocVersion{1.0}
\setDocStatus{Technical education / proposed reference}
\setDocOwner{CI/CD Engineering}
\setDocDate{30 October 2026}
\setDocSummary{Supplementary engineering practice}
\begin{document}
\makedoctitle
\makedocsummary
\tableofcontents
\clearpage
\section{Purpose, scope and authority}
This is an original modern engineering reference. The 2007 Duvall, Matyas and Glover book informs foundational CI intent but does not prescribe GitHub-specific implementations or 2026 supply-chain standards.

\section{Normative vocabulary}
Define change as a traceable revision; integration as verification of combined work; build as a script-defined transformation and checks; artifact as a versioned output; gate as a policy-enforced release condition. Distinguish triggering a workflow from achieving a successful integration.

Verification: collect the configuration or runtime evidence supporting this control, identify an accountable owner, and record the exception path where the condition cannot presently be met.

\section{Reference integration lifecycle}
Developer checks out a known revision, applies a small change, runs local fast validation and submits a pull request. A protected CI runner checks out that revision with minimal token permission, runs deterministic stages, records evidence and presents a decision. The team promptly repairs red mainline checks.

Verification: collect the configuration or runtime evidence supporting this control, identify an accountable owner, and record the exception path where the condition cannot presently be met.

\section{Design invariants}
Exactly one revision must identify an evaluated change. Every required gate must have an explicit pass/fail policy. Builds must not read undeclared secrets. Artifacts must be traceable to the successful evaluation. Re-runs do not erase prior failures; they create new evidence.

Verification: collect the configuration or runtime evidence supporting this control, identify an accountable owner, and record the exception path where the condition cannot presently be met.

\section{Failure-response protocol}
The first responder classifies failing tests, build infrastructure, dependency service or security gate. Record first-failed revision, owning team, timestamp, reproducer, mitigation and correction. Restore mainline health before layering additional unrelated changes.

Verification: collect the configuration or runtime evidence supporting this control, identify an accountable owner, and record the exception path where the condition cannot presently be met.

\section{Decision questions}
Can an engineer reproduce the pipeline locally? Are protected-branch checks current? How is a flaky test handled? Does the workflow exercise the intended revision? Which artifacts and logs survive a cancelled run?

Verification: collect the configuration or runtime evidence supporting this control, identify an accountable owner, and record the exception path where the condition cannot presently be met.

\section{Adoption checklist}
\begin{itemize}
\item Identify applicable repositories, application owners and trust boundaries.
\item Define minimum pass criteria and measurable acceptance thresholds.
\item Review workflow permissions, external actions, outputs and retention.
\item Record exceptions with justification, owner and expiry.
\item Revalidate after action, runner or policy changes.
\end{itemize}
\end{document}
